\section{Results: ordered Bi\texorpdfstring{$_2$}{2}Se\texorpdfstring{$_3$}{3} and Bi\texorpdfstring{$_2$}{2}Te\texorpdfstring{$_3$}{3} films}
\label{sec:results_ordered}

Ordered Bi$_2$Se$_3$ and Bi$_2$Te$_3$ films grown by TEM were used to validate the model by showing that the same detector-space intensity forward model quantitatively reproduces measured line shapes and intensities. Three detector images, collected at incident angles of $5^\circ$, $10^\circ$, and $15^\circ$, were used in the refinement. The $5^\circ$ condition gives the clearest single-view display of both the off-specular arc families and the specular $m=0$ features. Together, the data test the two-dimensional powder orientational limit, sample geometry, finite beam phase space, mosaicity, detector resolution, and ordered structure-factor weighting.

The detector-space panels and profile panels in Fig.~\ref{fig:ordered_chalcogenide_validation} use the same extraction convention. In panels~(a) and~(b), the solid colored curves are the calculated centerlines of selected $m$-indexed Bragg-rod trajectories on the detector. The semi-transparent colored bands are the finite detector regions integrated to form the one-dimensional profiles. The dashed colored outlines mark the edges of those integration bands. The central blue band is the $m=0$ specular-family trajectory, while the paired off-specular bands are the two detector-side intersections of the same two-dimensional powder rod family. The plus and minus labels identify those two detector-side branches only.

Panels~(c) and~(d) show the corresponding line profiles after applying the same projection and integration operation to the measured and calculated images. The horizontal coordinate is $L$, so a marker labeled $L=n$ denotes the nominal $00n$ or corresponding off-specular Bragg condition for that selected rod family. The black solid curves are the measured profiles. The orange dashed curves are the simulated profiles after the same detector-space projection. The off-specular rows are plotted on a linear intensity scale, while the $m=0$ row is plotted on a logarithmic scale because the important result is the persistence of weak, unexpected $00L$ intensity far from the strongest specular peak.

\begin{figure}[H]
  \centering
  % Consolidated ordered-film validation figure. Top row: detector-space
  % projection geometry. Bottom row: corresponding Q_z profile overlays.
  \begin{minipage}[t]{0.48\linewidth}
    \centering
    \textbf{Bi$_2$Se$_3$}\\[0.25em]
    \paneltagged{(a)}{\draftfigureplaceholder[0.32\textheight]{figures/results_ordered/figure7_bi2se3_qr_rod_qz_profiles_detector_selected_q_regions_5deg.png}{%
      Bi$_2$Se$_3$ detector image with projected fixed-$Q_R$ Bragg-rod centerlines and direct $HK$ labels. Use linear intensity, no stars, no legend, and absolute detector-pixel coordinates.
    }}
  \end{minipage}\hfill
  \begin{minipage}[t]{0.48\linewidth}
    \centering
    \textbf{Bi$_2$Te$_3$}\\[0.25em]
    \paneltagged{(b)}{\draftfigureplaceholder[0.32\textheight]{figures/results_ordered/figure7_bi2te3_qr_rod_qz_profiles_detector_selected_q_regions_5deg.png}{%
      Bi$_2$Te$_3$ detector image with projected fixed-$Q_R$ Bragg-rod centerlines and direct $HK$ labels. Use the same styling and label convention as Bi$_2$Se$_3$.
    }}
  \end{minipage}

  \vspace{0.8em}

  \begin{minipage}[t]{0.48\linewidth}
    \centering
    \paneltagged{(c)}{\draftfigureplaceholder[0.32\textheight]{figures/results_ordered/figure7_bi2se3_qr_rod_qz_profiles.png}{%
      Bi$_2$Se$_3$ measured/calculated $Q_z$ projection overlays for selected $HK$ families. Panels should use $Q_z$ or $L$ on the horizontal axis and matched measured/calculated curves.
    }}
  \end{minipage}\hfill
  \begin{minipage}[t]{0.48\linewidth}
    \centering
    \paneltagged{(d)}{\draftfigureplaceholder[0.32\textheight]{figures/results_ordered/figure7_bi2te3_qr_rod_qz_profiles.png}{%
      Bi$_2$Te$_3$ measured/calculated $Q_z$ projection overlays for selected $HK$ families. Match the Bi$_2$Se$_3$ axis, color, and labeling convention.
    }}
  \end{minipage}
  \caption[Ordered-film validation for Bi2Se3 and Bi2Te3]{Ordered-film validation of Bi$_2$Se$_3$ and Bi$_2$Te$_3$ at $\theta_i=5^\circ$. (a,b) Measured detector images with selected reciprocal-space extraction regions overlaid. Solid colored curves are the trajectory centerlines. Semi-transparent colored regions are the detector pixels integrated to form the profiles. Dashed colored outlines mark the integration-band boundaries. The labels $m$ identify two-dimensional powder rod families, and the signs label the two detector-side branches of the same family. The bright low-$L$ intensity near the critical-angle region is assigned to reflectivity caught by finite bandwidth. (c,d) Extracted profiles plotted versus $L$ after applying the same projection to measured and calculated images. Black solid curves are measured data and orange dashed curves are simulations. The off-specular rows test detector geometry, the two-dimensional powder construction, the mosaic-core width, and the ordered structure-factor pattern. The semilogarithmic $m=0$ rows emphasize the unexpected persistence of additional $00L$ peaks, which requires a long Lorentzian mosaic tail in addition to the narrow mosaic core.}
  \label{fig:ordered_chalcogenide_validation}
\end{figure}

Because mosaicity, beam divergence, wavelength spread, incidence-angle uncertainty, detector resolution, and finite integration width all broaden the effective scattering condition, the projected coordinate should not be interpreted as an ideal reciprocal-space cut from a single perfect crystallite. The comparison remains meaningful because the calculated image is subjected to the same projection and resolution effects before it is compared with the measured profile. The figure therefore tests the detector-space calculation as used experimentally. This distinction is especially important for the unexpected $00L$ peaks and the enhanced low-$L$ specular intensity, where incidence angle, wavelength bandwidth, divergence, detector-coordinate uncertainty, background, and the Lorentzian mosaic tail are most strongly coupled.

In Fig.~\ref{fig:ordered_chalcogenide_validation}, the off-specular and $m=0$ profiles play complementary roles. The off-specular arcs check the two-dimensional powder reciprocal-space construction and calibrated detector geometry, while their widths and asymmetries constrain the Gaussian-like mosaic core and instrumental resolution. The off-condition $m=0$ peaks test the Lorentzian mosaic tail and show why a rapidly decaying mosaic core alone is insufficient. The ordered-film result is therefore not only that the fitted curves match the strongest peaks, but that the same physical model also accounts for weaker and less intuitive detector features that were not built into a one-dimensional reduction.
